% paper/sections/sec4b_deep_unfolding.tex -- Balraj Singh (24BCE0531) -- runbook \S4.2 paper ownership
% \input this from paper/main.tex as Section 4's unfolding subsection (\subsection level).
% All numbers are macros populated by scripts/make_all_figures.py from results/ -- nothing
% hand-typed. \providecommand fallbacks exist only so the file compiles pre-integration.

\subsection{Deep Unfolding: A Learned Projected-Gradient Optimizer}

The model-free agents evaluated in Section~\ref{sec:results-drl} must discover from
scratch that maximizing the sum rate calls for ascent on
\eqref{eq:sum_rate}. Model-driven design instead builds that knowledge in
explicitly: we \emph{unfold} $L$ projected-gradient ascent steps on the true
objective into a fixed-depth network and train the few remaining degrees of
freedom end-to-end. Concretely, the $L$-layer optimizer executes
\begin{equation}
\boldsymbol{\theta}^{(l+1)}
\;=\; \operatorname{wrap}\!\Bigl(\boldsymbol{\theta}^{(l)}
+ \alpha_l \, \nabla_{\!\boldsymbol{\theta}}
R\bigl(\boldsymbol{\theta}^{(l)}\bigr)\Bigr),
\qquad
\alpha_l = e^{\rho_l} > 0,
\qquad l = 0, \dots, L-1,
\label{eq:unrolled_pg}
\end{equation}
with $L = 6$ layers, initialization
$\boldsymbol{\theta}^{(0)} = \boldsymbol{0}$ (identical to the AO bound of
Section~\ref{sec:ao_bound}, so the two are never separated by initialization
luck), and $\operatorname{wrap}(\boldsymbol{\theta}) =
\boldsymbol{\theta} - 2\pi\lfloor \boldsymbol{\theta}/2\pi \rfloor$, the
projection back onto $[0,2\pi)^N$ under which the rate
\eqref{eq:sum_rate} is exactly periodic. The only trainable parameters are the
$L$ real numbers $\rho_l$; the layer map itself is the exact gradient of the
frozen system model, evaluated by automatic differentiation through
\eqref{eq:zf}--\eqref{eq:sum_rate} with second-order graphs enabled so that the
training signal flows \emph{through} the internal gradients.

\subsubsection{Training Objective}
The network is trained to maximize the achievable rate itself:
\begin{equation}
\mathcal{L}_{\mathrm{rate}}
\;=\;
-\,\frac{1}{|\mathcal{B}|}\sum_{\boldsymbol{h} \in \mathcal{B}}
R\bigl(\boldsymbol{\theta}^{(L)}(\boldsymbol{h})\bigr),
\label{eq:rate_loss}
\end{equation}
minimized by Adam over mini-batches $\mathcal{B}$ of the $1500$ training
channels. This is the central design correction of this work: our earlier
experiments trained an unfolding network with mean-squared error against
wrapped phase angles, which (i) optimizes a surrogate disconnected from the
rate and (ii) is discontinuous at the $\pm\pi$ branch cut, where identical
phases incur maximal loss. We retain both defective objectives as ablations:
a smooth circular surrogate
$\mathcal{L}_{\mathrm{circ}} = \operatorname{mean}\bigl(1 -
\cos\bigl(\operatorname{wrap}(\boldsymbol{\theta}^{(L)} -
\boldsymbol{\theta}_{\mathrm{AO}})\bigr)\bigr)$ against the AO reference
phases, and the original wrapped-angle MSE
$\mathcal{L}_{\mathrm{MSE}} = \operatorname{mean}\bigl(\|
\operatorname{wrap}\,\boldsymbol{\theta}^{(L)} -
\operatorname{wrap}\,\boldsymbol{\theta}_{\mathrm{AO}} \|^2\bigr)$,
reproducing its deficit in Section~\ref{sec:results_bound}.

\subsubsection{Why Unfolding Decides Fast}
At inference the unfolded optimizer performs exactly $L = 6$ fixed
gradient evaluations---no line search, no iterative-until-convergence
loop---on a single frozen test channel, and its latency
is therefore a deterministic constant independent of channel difficulty. The
Riemannian AO bound of Section~\ref{sec:ao_bound}, by contrast, iterates with
Armijo line searches until a first-order tolerance is met. The measured
per-decision latencies and the resulting quality--latency trade are reported
in Section~\ref{sec:results_bound}; the claim we defend is the \emph{relative}
one (unfolding is orders of magnitude faster than the bound it approximates),
with the absolute milliseconds table reported as measured on the training
dtype.

\subsubsection{Protocol}
The unfolding is trained on the $1500$ training channels with the identical
budget, evaluation schedule, frozen test set, and five-seed protocol as every
model-free agent, and the same one-hyperparameter tuning rule (the Adam
learning rate over a three-point grid). The circular and MSE ablations
additionally require AO reference phases on the training split, which are
produced once by the reference script and stored under \texttt{results/ao/}.

% ---- auto-populated macros (make_all_figures.py) ------------------------------
\providecommand{\UnfoldingLayers}{6}
\providecommand{\UnfoldLRgrid}{$\{3\cdot10^{-2},\,10^{-2},\,3\cdot10^{-3}\}$}
The depth is $L = \UnfoldingLayers$ and the learning-rate grid is
\UnfoldLRgrid{}; no other hyperparameter is tuned.
